\subsection{Orthonormalisation}

THEOREM 2. --- For every orthonormal set L in a hilbertian space E, there exists an orthonormal basis B of E containing L.

In fact, let D be the family of all orthonormal subsets of E, linearly ordered by inclusion; it is immediate that this family has finite character (S, III, § 4, No.~5). Hence there exists a maximal family B in D containing L, by th. 1 of S, III, § 4, No.~5. It remains to prove that B is a total set. If not, there will exist a vector \(y \neq 0\) which is orthogonal to all the vectors of B (V, p.~22, prop. 5), and multiplying y by a suitable scalar, we can assume that \(\|y\| = 1\) ; then, \(B \cup \{y\}\) will be an orthonormal set distinct from B and containing B; this contradicts the definition of B; hence the theorem.

COROLLARY 1. --- In every hilbertian space, there exists an orthonormal basis.

It is enough to apply th. 2 to the case L = ∅.

COROLLARY 2. --- Every hilbertian space is isomorphic to a space \(\ell^2 (\mathbf{I})\) .

More precisely, let \((e_{i})_{i\in I}\) be an orthonormal basis of a hilbertian space E. By props. 4 (V, p.~21) and 5 (V, p.~22), the mapping \(\phi\) defined by

\[
\phi (x) = \big (\langle e _ {i} | x \rangle \big) _ {i \in \mathrm{I}}\tag{4}
\]

is a hilbertian space isomorphism from E onto \(\ell^2 (\mathrm{I})\) . The inverse isomorphism \(\psi\) is defined by

\[
\psi ((\lambda_ {i}) _ {i \in \mathbf {I}}) = \sum_ {i \in \mathbf {I}} \lambda_ {i} e _ {i}.\tag{5}
\]

PROPOSITION 6. --- Let E be a Hausdorff prehilbertian space, and let \((a_{n})_{n\in\mathbb{I}}\) (I an interval of N with origin 1) be a countable (finite or not) independent family of vectors of E. There exists an orthonormal family \((e_{n})_{n\in\mathbb{I}}\) , and only one, in E, with the following properties :

\begin{enumerate}
\def\labelenumi{\arabic{enumi})}
\item
  for every integer \(p \in \mathbf{I}\) , the vector subspace of \(\mathbf{E}\) generated by \(e_1, e_2, \ldots, e_p\) is identical with the vector subspace of \(\mathbf{E}\) generated by \(a_1, a_2, \ldots, a_p\) ;
\item
  for every integer \(p \in \mathbf{I}\) , the number \(\langle a_p | e_p \rangle\) is real and \(> 0\) .
\end{enumerate}

In fact, let \(V_{n}\) be the subspace (of dimension n) generated by \(a_{1}, a_{2}, ..., a_{n}\) . If \(n + 1 \in I\) and \(b_{n+1} = a_{n+1} - p_{\mathrm{V}_{n}}(a_{n+1})\) (where \(p_{V_{n}}\) is the orthoprojector onto the complete subspace \(V_{n}\) ), the line \(Kb_{n+1}\) is the orthogonal of \(V_{n}\) in \(V_{n+1}\) . If the \(e_{n}\) satisfy condition 1) of the proposition, we must have \(e_{n+1} = \lambda b_{n+1}\) ; the condition \(\|e_{n+1}\| = 1\) then implies \(|\lambda|^2 \|b_{n+1}\|^2 = 1\) and the condition \(\langle a_{n+1}|e_{n+1} \rangle > 0\) implies \(\lambda \langle a_{n+1}|b_{n+1} \rangle > 0\) ; this completely determines \(\lambda\) , and we have proved that we can determine, by induction, an orthonormal family \((e_n)_{n\in I}\) , and only one, so as to satisfy conditions 1) and 2) of the proposition.

The sequence \((e_{n})_{n\in\mathbb{I}}\) is said to be obtained by orthonormalisation from the independent family \((a_{n})_{n\in\mathbb{I}}\) . It is clear that the vector subspace generated by the family \((e_{n})\) is identical with the vector subspace generated by the family \((a_{n})\) . In particular, if \((a_{n})\) is a total sequence, so is \((e_{n})\) , which is then an orthonormal basis of E; hence we get:

COROLLARY. --- In every Hausdorff prehilbertian space E satisfying the first axiom of countability, there exists a countable orthonormal basis.

If E satisfies the first axiom of countability, then there exists a total sequence in E, and we can always extract an independent total family from such a sequence (A, II, § 7, No.~1, th. 2).

We can give examples of Hausdorff prehilbertian spaces not having any orthonormal basis (V, p.~70, exerc. 2).

Example. --- Let I be the interval \((-1,1)\) of R and E the vector space of real valued continuous functions on I. Let x denote the canonical injection from I into R, considered as an element of E. By the Stone-Weierstrass theorem the sequence \((x^{n})_{n\in\mathbf{N}}\) is total in E for the topology of uniform convergence GT, X, § 4, No.~2).

Consider E as a real Hausdorff prehilbertian space in which the scalar product is given by

\[
\langle f | g \rangle = \int_ {- 1} ^ {1} f (t) g (t) d t.
\]

The sequence \((x^{n})_{n\in \mathbb{N}}\) is then total in the prehilbertian space E. Let \((\Pi_n)_{n\in \mathbb{N}}\) be the sequence obtained by the orthonormalisation of the sequence \((x^n)_{n\in \mathbb{N}}\) . We can show that \(\Pi_n = (n + \frac{1}{2})^{1 / 2}\mathbf{P}_n\) , where the Legendre polynomial \(\mathbf{P}_n\) is defined by

\[
\mathrm{P} _ {n} (x) = \frac {1}{2 ^ {n} n !} \left(\frac {d}{d x}\right) ^ {n} (x ^ {2} - 1) ^ {n}.
\]

PROPOSITION 7. --- In a hilbertian space E, two orthonormal bases are equipotent.

Let B and C be two orthonormal bases of E. The case when one of the two sets B, C is finite is trivial, since a finite orthonormal basis is an algebraic basis of the space. Suppose therefore that B and C are infinite. For every \(x \in B\) , let \(C_{x}\) be the subset of C consisting of all \(y \in C\) such that \(\langle x | y \rangle \neq 0\) . The set \(C_{x}\) is countable (V, p.~21, prop. 4). For every \(y \in C\) , there exists \(x \in B\) such that \(y \in C_{x}\) , since B is an orthonormal basis and \(y \neq 0\) ; in other words C is the union of the countable sets \(C_{x}\) as x ranges over B. The cardinality of C is hence less than that of N × B, hence less than that of B (S, III, §6, No.~3, cor. 4); similarly, the cardinality of B is less than that of C; this completes the proof.

The cardinality of an arbitrary orthonormal basis of a hilbertian space E is called the hilbertian dimension of E.

COROLLARY 1. --- Given two orthonormal bases in a hilbertian space E, there exists an automorphism of E transforming the first basis into the second.

COROLLARY 2. --- In order that the hilbertian spaces \(\ell^2 (\mathbf{I})\) and \(\ell^2 (\mathbf{J})\) be isomorphic, it is necessary and sufficient that I and J are equipotent.

